
\documentclass[10pt,a4paper]{extarticle}
\usepackage{parskip}

% ======================
% Pacotes e configuração
% ======================
\usepackage[utf8]{inputenc}
\usepackage[T1]{fontenc}
\usepackage[brazil]{babel}
\usepackage{amsmath, amssymb, amsthm}
\usepackage{geometry}
\usepackage{enumitem}
\usepackage{booktabs}
\usepackage{hyperref}
\geometry{margin=2.5cm}
\setlist{nosep}

\begin{document}

\begin{center}
{\Large \textbf{Estudo Dirigido — Condições de Otimalida para Problemas sem Restrições}}\\[4pt]
\textit{Classificação de Matrizes, Formas Quadráticas e Condições de Segunda Ordem}\\[6pt]
\end{center}

\noindent\textbf{Observação:} Você pode consultar tudo o que estiver escrito de próprio punho no caderno da disciplina. 

\bigskip

\begin{enumerate}[label=\textbf{Q\arabic*.}]

\item Seja \( f:\mathbb{R}^2 \to \mathbb{R} \) uma função contínua e diferenciável, com domínio \( D \subseteq \mathbb{R}^2 \).Defina formalmente o que significa dizer que \( x^\star \in D \) é um \textbf{mínimo local}, um \textbf{máximo local}, um \textbf{mínimo global} e um \textbf{máximo global} de \( f \).

\item Utilize as definições de mínimo e máximo global para mostrar que \texttt{min }\(f(x) \), \( f:\mathbb{R}^2 \to \mathbb{R} \) equivale a \texttt{max }\(1/f(x)\), se possível.

% Q1
\item Considere a matriz simétrica
\[
A = \begin{bmatrix}
4 & -2 & 0\\
-2 & 3 & -1\\
0 & -1 & 2
\end{bmatrix}.
\]
\begin{enumerate}[label=(\alph*)]
\item Calcule o polinômio característico $\chi_A(\lambda)=\det(A-\lambda I)$ e seus autovalores.
\item Classifique $A$ quanto à definitude (PD, PSD, ND, NSD, indefinida), com base no sinal dos autovalores.
\item Escreva a forma quadrática $q(x)=x^\top A x$ e discuta o sinal de $q(x)$ para $x\neq 0$ à luz da sua classificação.
\end{enumerate}

\bigskip

% Q2
\item Seja
\[
B(t)=\begin{bmatrix}
2 & t & 0\\
t & 3 & 1\\
0 & 1 & 1
\end{bmatrix},\qquad t\in\mathbb{R}.
\]
Usando o critério de Sylvester (menores principais líderes), determine os valores de $t$ para os quais $B(t)$ é:
\begin{enumerate}[label=(\alph*)]
\item definida positiva (PD);
\item semidefinida positiva (PSD);
\item indefinida.
\end{enumerate}

\bigskip

% Q4
\item Sejam
\[
f(x,y)= x^4 + y^4 - 4x^2 + 2xy + 1.
\]
\begin{enumerate}[label=(\alph*)]
\item Encontre todos os pontos críticos resolvendo $\nabla f(x,y)=0$.
\item Compute a Hessiana $\nabla^2 f(x,y)$ e, em cada ponto crítico, classifique-o (mínimo local estrito, máximo local estrito ou ponto de sela) usando a definitude da Hessiana.
\item Explique, com base nas notas teóricas, por que a definitude da Hessiana em $x^\star$ determina a natureza do ponto crítico.
\end{enumerate}

\bigskip

% Q5
\item Seja $A\in\mathbb{R}^{n\times n}$ simétrica. Prove \emph{ou esboce uma prova rigorosa} da equivalência abaixo:

\texttt{$A$ é definida positiva $\iff$ todos os autovalores de $A$ são estritamente positivos.}

\item Considere uma função \( f:\mathbb{R}^2 \to \mathbb{R} \) cuja análise em um ponto \( x^\star = (1, -1) \) revelou:
\[
\nabla f(x^\star) =
\begin{bmatrix}
0\\
0
\end{bmatrix},
\qquad
\nabla^2 f(x^\star) =
\begin{bmatrix}
2 & -1\\
-1 & 3
\end{bmatrix}.
\]
\begin{enumerate}[label=(\alph*)]
\item O que podemos concluir sobre o ponto \( x^\star \)?
\item Justifique com base nos autovalores da Hessiana.
\item Interprete o comportamento local de \( f \) em torno de \( x^\star \) utilizando uma expansão em série de Taylor.
\end{enumerate}

\item Considere uma função \( f:\mathbb{R}^2 \to \mathbb{R} \) cuja análise em um ponto \( x^\star = (0, -1) \) revelou:
\[
\nabla f(x^\star) =
\begin{bmatrix}
-0.5\\
0.1
\end{bmatrix},
\qquad
\nabla^2 f(x^\star) =
\begin{bmatrix}
1 & 0\\
0 & 1
\end{bmatrix}.
\]
\begin{enumerate}[label=(\alph*)]
\item O que podemos concluir sobre o ponto \( x^\star \)?
\item Interprete o comportamento local de \( f \) em torno de \( x^\star \) utilizando uma expansão em série de Taylor.
\end{enumerate}

\end{enumerate}


\end{document}
